\hspace{3mm}

\centerline{\bf \Huge Домашнее задание.}
\centerline{\bf \Huge Переменные и логика.}

\begin{flushright}
    \scriptsize{Если я сдамся из-за того, что испугался,\\
                то не смогу стать сильнее.
                \\Сон Джин-Ву. Поднятие уровня в одиночку}
\end{flushright}

\problem{1!} Известно, что все шиншиллы говорят только правду. Верно ли утверждение: «Все, кто говорят правду, являются шиншиллами»?
Объясните, почему?

\problem{2!} У АА обучаются 42 ученика(шиншиллы и лососи). Известно, что шиншилл в два раза больше, чем лососей. Сколько шиншилл обучаются у АА?

\problem{3} В ЭлМШ обучаются 48 учеников, причем каждый принадлежит только своей группе. Каждый второй ученик является шиншиллой, каждый третий - лососем, а каждый четвертый - киви. Сколько шиншилл, сколько лососей и сколько киви обучаются в ЭлМШ?

\problem{4} В ЭлМШ шиншиллы всегда говорят правду, а лососи всегда лгут. На встрече ЭлМШ ровно половина учеников сказали, что среди всех, кто пришел, нет лжецов. Сколько лососей и сколько шиншилл пришло на встречу ЭлМШ, если всего на встрече было 50 учеников? Разберите все возможные варианты!!!